
We asked whether the reported marginality of in-season replanning survives
when the biological outcome responds sharply to timing error. Across
10 independently generated farms it does not: the gain from adaptive
replanning rises with the curvature of the yield response on every farm
($r = 0.94$) and reaches five percent at roughly 1.8 times the
curvature measured for hybrid maize. The irrigation finding is therefore
bounded rather than general, and we locate the boundary.

For multi-field hybrid maize seed production at its measured response,
adaptive replanning is worth 1.91\% (SD 1.49, $p = 0.003$) against a
robust static plan --- positive on 9 of 10 farms, significant
on 5, and largest where the static plan performs worst. Below the
boundary, a well-built static plan is the better engineering choice.

This is the first formulation of hybrid seed sowing as a resource-constrained
stochastic scheduling problem coupled to parent flowering synchrony, scored
against a published field regression and solved to proven optimality so that
the measured difference is attributable to adaptivity rather than to search.

A methodological point stands alongside the biological one. Measured on a
single farm, the same system showed an effect roughly twice its size and a
saturation result that did not replicate. Instance-level replication decided
both.
